% Generated draft: requires mathematical and source review.
\documentclass[11pt]{article}
\usepackage[margin=25mm]{geometry}
\usepackage{fontspec}
\setmainfont{Latin Modern Roman}
\usepackage{amsmath,amssymb,mathtools,unicode-math}
\setmathfont{Latin Modern Math}
\usepackage{enumitem,needspace,xurl,hyperref}
\setlength{\parindent}{0pt}
\setlength{\parskip}{5pt}
\setlength{\emergencystretch}{3em}
\newcommand{\N}{\mathbb{N}}
\newcommand{\Z}{\mathbb{Z}}
\newcommand{\Q}{\mathbb{Q}}
\newcommand{\R}{\mathbb{R}}
\newcommand{\C}{\mathbb{C}}
\newcommand{\F}{\mathbb{F}}
\newcommand{\A}{\mathbb{A}}
\newcommand{\PP}{\mathbb{P}}
\newcommand{\E}{\mathbb{E}}
\newcommand{\eps}{\varepsilon}
\newcommand{\dd}{\,\mathrm{d}}
\begin{document}
\section*{UnsolvedMath ID: 11000105}
\textbf{Source problem code:} catalogue entry
\subsection{Definitions and mathematical statement}
Let \(G=\operatorname{Mod}_S\) be the mapping class group of the surface \(S\) intended by the source. For each positive integer \(N\), define the power subgroup
\[
 P_N(G)=\langle x^N:x\in G\rangle.
\]
It is characteristic, in particular normal. Every element of \(G/P_N(G)\) has order dividing \(N\).

\emph{Question I.} Does there exist \(N_0\) such that
\[
 [G:P_N(G)]=\infty\qquad\text{for every integer }N\geq N_0?
\]
This is the usual quantified reading of ``for sufficiently big \(N\).'' The threshold may depend on \(S\). The statement does not merely ask for one exponent, an unbounded subsequence, or sufficiently divisible exponents; those would be distinct weakenings.

As the source notes, an affirmative infinite-index instance gives an infinite quotient \(G/P_N(G)\) of bounded exponent dividing \(N\), described there as an infinite Burnside group. It is not thereby asserted to be the free Burnside group, or to contain elements of every order dividing \(N\). The type and complexity of \(S\), boundary/marking conventions, and any excluded low-complexity cases are absent; these hypotheses are essential to determine the intended group and the validity of a uniform-large-exponent statement. The leading extracted ``s'' is read as the truncated beginning of ``Is.
\subsection{Short English statement}
For all sufficiently large exponents, is the subgroup generated by all \(N\)-th powers in the mapping class group of infinite index? The corresponding quotient would be an infinite group of exponent dividing \(N\); the surface hypotheses are missing.
\subsection{Sources}
\begin{itemize}
\item Local UnsolvedMath registry, permanent ID 11000105, source code catalogue entry, Question I, PDF page 87. No individual proposer is supplied. Page identifies the item despite repeated labels across chapters.
\item Farb, \emph{Problems on Mapping Class Groups and Related Topics}. Supplied URL: \url{https://www.math.uchicago.edu/~farb/papers/mcgbook.pdf}; recorded access 2026-07-29; extraction pdf-font-structure. The edited volume contains 329 explicitly typeset headings across contributed chapters.
\item Supplied status open and OPEN-TRIAGE. The research summary mentions power subgroups and says open, without a precise reference or theorem. Current primary-literature status and redistribution terms for the publicly accessible author/publisher copy remain NEEDS\_REVIEW.
\end{itemize}
\end{document}
